\documentclass[11pt,a4paper]{article}
\usepackage[utf8]{inputenc}
\usepackage[T1]{fontenc}
\usepackage{amsmath,amssymb,amsfonts}
\usepackage{booktabs}
\usepackage{graphicx}
\usepackage{hyperref}
\usepackage{geometry}
\geometry{margin=1in}

\title{\textbf{Traceable Hybrid Auditing of Clinical LLM Outputs: A Proof-of-Concept Combining Span Verification, Protocol Rules, and Graph Reasoning Kernels}}

\author{
  \textbf{Hammam Omer} \\
  Faculty of Medicine, Omdurman Islamic University \\
  Omdurman, Sudan \\
  \texttt{Repository: \url{https://github.com/HammamNM/reasoning-debugger}}
}
\date{October 2026}

\begin{document}

\maketitle

\begin{abstract}
\noindent \textbf{Background:} Evaluating Large Language Models (LLMs) in open-ended clinical workflows such as acute chest pain triage requires auditing procedural adherence (omissions, temporal ordering, and contraindications) rather than static multiple-choice accuracy. Unconstrained \textit{LLM-as-a-Judge} paradigms risk fabricated event extraction and often penalize appropriate diagnostic uncertainty. \\
\textbf{Objective:} We present a proof-of-concept \textbf{Traceable Neuro-Symbolic Auditing Framework} and evaluate its mechanism calibration on a pilot benchmark ($N=10$ traces, plus $N=5$ linguistic edge cases). \\
\textbf{Methods:} The architecture decouples auditing into three layers: (1) a \texttt{TraceExtractor} enforcing deterministic verbatim span verification and epistemic assertion calibration (\texttt{PROPOSED}, \texttt{REPORTED\_PERFORMED}, \texttt{NEGATED}, \texttt{SUSPECTED}, \texttt{DEFINITIVE}), (2) a \texttt{ClinicalRuleEvaluator} checking temporal ordering and contraindications, and (3) a graph-based \texttt{Reasoning Kernel} detecting unsupported definitive claims (\texttt{Pattern RP-001}). \\
\textbf{Pilot Results:} Across our 10-case pilot calibration suite (5 synthetic failure/control traces and 5 real \texttt{GPT-4o} responses), introducing assertion-status calibration eliminated 3 false-positive flags on negated actions and suspected differentials, shifting pilot specificity from 4/7 (57.1\%, 95\% CI: 18.4\%--90.1\%) to 7/7 (100.0\%, 95\% CI: 59.0\%--100.0\%) while preserving 3/3 sensitivity on synthetic protocol violations.
\end{abstract}

\section{Introduction}
While Large Language Models (LLMs) perform strongly on static medical licensing examinations \cite{kung2023performance}, open-ended emergency triage of \textbf{Acute Chest Pain} \cite{gulati2021aha} requires strict procedural adherence. An LLM may list a correct differential diagnosis while proposing lethal management steps, such as discharging a patient after symptomatic relief from a gastrointestinal cocktail prior to reviewing serial troponins (premature closure), or administering empiric heparin in suspected aortic dissection.

Current \textit{LLM-as-a-Judge} evaluation methods \cite{zheng2023judging} face two structural limitations in procedural clinical auditing:
\begin{enumerate}
    \item \textbf{Unverified Event Extraction:} A neural judge may hallucinate that a candidate response ordered an investigation when no such statement exists in the text.
    \item \textbf{Epistemic Conflation:} Inspired by classical clinical negation and context detection algorithms such as \texttt{NegEx} \cite{chapman2001simple} and \texttt{ConText} \cite{harkema2009context}, an effective auditor must distinguish negated warnings (\textit{``Do not give heparin''}) and suspected differentials from definitive clinical commitments.
\end{enumerate}

\section{System Architecture}
Our framework separates linguistic extraction from deterministic symbolic verification across three tiers:

\subsection{Tier 1: \texttt{TraceExtractor} (Span Firewall \& Epistemic Tagging)}
Given raw text $T$, candidate events $E = \{e_1, \dots, e_n\}$ must satisfy verbatim substring verification under whitespace/case normalization $\mathcal{N}(s_i) \subseteq \mathcal{N}(T)$, rejecting ungrounded extractions (\texttt{SPAN\_NOT\_IN\_SOURCE}). Lexical synonyms are mapped to canonical protocol codes, and events are filtered by assertion status $\alpha(e_i)$ so that \texttt{NEGATED} warnings are never penalized as proposed actions, and only \texttt{DEFINITIVE} inferences require structural evidentiary proof.

\subsection{Tier 2 \& Tier 3: \texttt{ClinicalRuleEvaluator} \& \texttt{Reasoning Kernel}}
The rule engine evaluates mandatory omissions (\texttt{CRITICAL\_OMISSION}), prerequisite temporal ordering $(a_{\text{pre}} \prec a_{\text{target}})$ for premature closure, and contraindications (\texttt{must\_avoid}). Concurrently, the \texttt{Reasoning Kernel} constructs a Directed Acyclic Graph $G = (V, E_{\text{graph}})$ and triggers \texttt{RP-001} whenever a \texttt{DEFINITIVE} inference node has support in-degree $\deg^{-}_{\text{SUPPORTS}}(v) = 0$.

\section{Pilot Calibration \& Mechanism Verification}
Table \ref{tab:ablation} reports the pilot mechanism ablation across $N=10$ calibration traces (3 synthetic protocol violations, 2 synthetic safe traces, and 5 real \texttt{GPT-4o} responses) alongside an unseen 5-case linguistic stress suite (\texttt{run\_edge\_cases.py}).

\begin{table}[h!]
\centering
\caption{Pilot Mechanism Ablation: Assertion Calibration on $N=10$ Traces}
\vspace{0.2cm}
\resizebox{\textwidth}{!}{%
\begin{tabular}{lccccccc}
\toprule
\textbf{Auditor Configuration} & \textbf{TP} & \textbf{FP} & \textbf{TN} & \textbf{FN} & \textbf{Recall (3/3)} & \textbf{Specificity (TN/7)} & \textbf{Pilot $F_1$} \\
\midrule
1. Uncalibrated Baseline & 3 & 3 & 4 & 0 & 100.0\% (29.2--100\%) & 57.1\% (18.4--90.1\%) & 0.667 \\
\textbf{2. Calibrated Hybrid Pipeline} & \textbf{3} & \textbf{0} & \textbf{7} & \textbf{0} & \textbf{100.0\% (29.2--100\%)} & \textbf{100.0\% (59.0--100\%)} & \textbf{1.000} \\
\bottomrule
\end{tabular}%
}
\label{tab:ablation}
\end{table}

\section{Methodological Limitations \& Future Scaling}
This work represents a software architecture and pilot calibration study rather than a large-scale empirical validation. Specifically: (1) the sample size ($N=10$ calibration traces, $N=5$ edge cases) yields wide Clopper-Pearson 95\% confidence intervals; (2) rules and synthetic test cases were authored by a single investigator without independent multi-clinician inter-annotator agreement ($\kappa$); and (3) while substring verification confirms quote existence, upstream assertion classification remains dependent on parser accuracy. Future work will expand evaluation to $>50$ multi-disease scenarios with blinded dual-clinician annotation and direct comparison against standalone LLM judges.

\begin{thebibliography}{9}
\bibitem{gulati2021aha}
M. Gulati \textit{et al.}, ``2021 AHA/ACC Guideline for the Evaluation and Diagnosis of Chest Pain,'' \textit{Circulation}, vol. 144, no. 22, pp. e368--e454, 2021.
\bibitem{chapman2001simple}
W. W. Chapman \textit{et al.}, ``A simple algorithm for identifying negated findings and diseases in discharge summaries,'' \textit{Journal of Biomedical Informatics}, vol. 34, no. 5, pp. 301--310, 2001.
\bibitem{harkema2009context}
H. Harkema \textit{et al.}, ``ConText: An algorithm for determining negation, experiencer, and temporal status from clinical reports,'' \textit{Journal of Biomedical Informatics}, vol. 42, no. 5, pp. 839--851, 2009.
\bibitem{kung2023performance}
T. H. Kung \textit{et al.}, ``Performance of ChatGPT on USMLE,'' \textit{PLOS Digital Health}, vol. 2, no. 2, p. e0000198, 2023.
\bibitem{zheng2023judging}
L. Zheng \textit{et al.}, ``Judging LLM-as-a-Judge with MT-Bench and Chatbot Arena,'' in \textit{NeurIPS}, 2023.
\end{thebibliography}

\end{document}
